\section{Background and related work}\label{sec:related}

\paragraph{Personal-data detection and de-identification}
Rule-based de-identification of clinical text dates back to lookup tables and regular expressions applied to nursing notes \cite{neamatullah2008deid}; Philter refined the approach on a large annotated clinical corpus \cite{norgeot2020philter}. The 2014 i2b2/UTHealth shared task established annotated corpora and evaluation practice \cite{stubbs2015annotating}, and recurrent neural networks reached strong results on them \cite{dernoncourt2017deid}. Presidio \cite{presidio} combines pattern recognisers with spaCy or transformer models and ships an evaluation toolkit with a synthetic dataset \cite{presidioresearch}. Lison et al.\ argue that sequence labelling alone does not amount to anonymisation and call for explicit disclosure-risk measures \cite{lison2021anonymisation}; the Text Anonymization Benchmark operationalises this for court judgments \cite{pilan2022tab}. Named-entity recognition underpins the free-text part of these systems, from the CoNLL-2003 shared task \cite{tjongkimsang2003conll} to transformer encoders \cite{devlin2019bert,sanh2019distilbert} and generalist span models \cite{zaratiana2024gliner}. Pseudonymisation with a keyed hash such as HMAC \cite{rfc2104} preserves referential integrity but reveals equality between values, so it is not anonymisation in the sense of $k$-anonymity \cite{sweeney2002kanonymity}. Redaction as studied here is narrower than anonymisation: it removes spans that detectors recognise and makes no claim about re-identification from what remains.

\paragraph{Credential and secret detection}
Meli et al.\ measured secret leakage across public GitHub repositories \cite{meli2019git}. Approaches range from vendor-specific regular expressions \cite{sinha2015secretkey} to machine-learned filters that reduce false positives \cite{saha2020secrets} and neural detectors for free-form passwords \cite{feng2022passfinder}. SecretBench provides a labelled corpus of candidate secrets \cite{basak2023secretbench}, and a comparison of nine tools attributed false positives to generic patterns and entropy heuristics and false negatives to missing rules \cite{basak2023comparative}. Krause et al.\ studied how developers handle leaked secrets \cite{krause2023pushed}. Shannon entropy \cite{shannon1948} is the usual signal for undocumented high-entropy secrets. The Nosey Parker rule set \cite{noseyparker} supplies the labelled examples we evaluate on.

\paragraph{Detection across segment boundaries}
That a pattern split across units of input can evade a detector is well known in network security: Ptacek and Newsham showed that intrusion detectors which do not reconstruct streams as the end host does can be evaded \cite{ptacek1998insertion}, and Hyperscan \cite{wang2019hyperscan} provides a streaming mode whose matches do not depend on how the data is split into blocks \cite{intel2019hyperscanintro}. These systems carry automaton state across boundaries, which makes pattern matching exact but says nothing about what happens after a match: validators, context-dependent confidences and joint overlap resolution all read text the automaton has already consumed. A redaction pipeline differs in three ways that affect correctness: detectors consult bounded context on both sides of a match, confidences depend on nearby text, and overlapping detections are resolved jointly before the output is rewritten. Overlapping windows are a common way to process long inputs, and the stream model \cite{muthukrishnan2005streams} formalises computation in bounded memory; linear-time regular-expression matching goes back to Thompson \cite{thompson1968regex}. Among the redaction, de-identification and secret-scanning systems cited above, we found none that states when chunked output equals whole-document output or how the overlap must be sized. We do not claim that the underlying locality argument is new; the contribution is its instantiation and enforcement for context-dependent redaction.
